\documentclass[10pt,a4paper]{amsart}
\usepackage[margin=19mm]{geometry}
\usepackage{amsmath,amssymb,mathtools}
\usepackage[scr=rsfs,cal=euler]{mathalfa}
\usepackage{xfrac}
\usepackage[unicode,hidelinks,bookmarks=false]{hyperref}
\newcommand{\ie}{i.e.\ }
\newcommand{\cf}{cf.\ }
\newcommand{\resp}{respectively\ }
\newcommand{\Subsection}{\subsection}
\providecommand{\nobreakdash}{\nobreak\hbox{-}}
\setlength{\emergencystretch}{2em}
\multlinegap0pt
\usepackage{bm}
\usepackage{bbm}
\usepackage{dsfont}
\usepackage{xspace}
\usepackage{cancel}
\usepackage{mathtools}
\usepackage{amsthm}
\makeatletter
\@ifundefined{theorem}{\newtheorem{theorem}{Theorem}}{}
\@ifundefined{lemma}{\newtheorem{lemma}[theorem]{Lemma}}{}
\@ifundefined{proposition}{\newtheorem{proposition}[theorem]{Proposition}}{}
\makeatother
\newcommand{\R}{\mathbb{R}}
\title[On dispersive decay for the generalized Korteweg--de Vries e]{On dispersive decay for the generalized Korteweg--de Vries equation}
\author{Matthew Kowalski, Minjie Shan}
\date{Source: arXiv:2510.01728v1 (CC BY 4.0)}
\begin{document}
\maketitle

\noindent\emph{Source and licence.} On dispersive decay for the generalized Korteweg--de Vries equation. Version arXiv:2510.01728v1 (CC BY 4.0), distributed under the Creative Commons Attribution 4.0 International licence, as recorded in the arXiv licence metadata for this exact version and shown on the abstract page https://arxiv.org/abs/2510.01728v1. Full rights notice: licenses/arxiv-2510.01728-CC-BY-4.0.txt.\par\smallskip

\noindent\textit{Editorial scope.} Counted unit: the Proposition whose source label is \texttt{prel/strichartz} in the article (source0.tex lines 333--345, source label \texttt{prel/strichartz}), delivered together with the article's own complete original proof of it (source0.tex lines 346--379, one \texttt{proof} environment of 34 lines). No mathematics is written, repaired or re-derived: statement and proof are the article's own text line for line. Required context retained uncounted: prel/strichartz-basic (lines 325--330). Every definition, display and label of those lines is retained unchanged. Omitted from this package and not redistributed: 1--324, 331--332, 380--552. Nothing inside a retained excerpt is omitted or altered: every retained body is delivered exactly as the article's lines are written, so no omission receipt of any kind accompanies this package. Citations: the retained excerpts carry no citation command at all, so no bibliography entry of the article is transcribed and no reference is redistributed. Reference adaptation: none was needed and none was made; every reference inside the delivered unit resolves inside it, so no unresolved reference or citation is produced. Printed slips kept literal: no printed word, symbol, hypothesis or number of the counted statement or of its proof was altered. Layout and class: the article's own class is \texttt{amsart} with packages including amsfonts, amsmath, amssymb, appendix, mathrsfs, hyperref, graphicx, amsthm, color, dsfont, bbm, cite; the wrapper is the lane's pinned \texttt{amsart} editorial wrapper at 10pt on A4, which restates the theorem-like environments the retained text needs and adds the article's own macro definitions that the retained text needs (\texttt{\textbackslash R}), with extra wrapper packages (bm, bbm, dsfont, xspace, cancel, mathtools). The delivered numbering is the unit's own. Assets: no retained span contains a figure environment, a graphics-inclusion command, a PDF-page inclusion, a table or a TikZ picture; no image file is copied into this package, the package declares no asset dependency and no image file is redistributed with it. Licence: version arXiv:2510.01728v1 of this article is distributed under the Creative Commons Attribution 4.0 International licence, as recorded in the arXiv licence metadata for this exact version and shown on the abstract page https://arxiv.org/abs/2510.01728v1. A full rights notice is delivered at licenses/arxiv-2510.01728-CC-BY-4.0.txt. Characters: every retained excerpt is pure ASCII, so the delivered engine, XeLaTeX with its default Latin Modern text font, reports no missing character.
\begin{lemma}\label{prel/strichartz-basic}
    Let $\theta \in [0,1)$, $\alpha \in [0,\frac{1}{2}]$, $(q,p) = (\frac{6}{\theta(\alpha + 1)},\frac{2}{1-\theta})$ and $\frac{1}{p} + \frac{1}{p'} = \frac{1}{q} + \frac{1}{q} = 1$. Then for any spacetime slab $J \times \R$ and $t_0 \in J$,
    \begin{equation*}
        \bigg\|\int_{t_0}^t |\partial_x|^\frac{\theta\alpha}{2} e^{-(t-s)\partial_x^3} g(s,\cdot) ds\bigg\|_{L_t^{q,2}L_x^{p}(J)} \lesssim \big\||\partial_x|^{-\frac{\theta\alpha}{2}}g\big\|_{L_t^{q',2}L_x^{p'}(J)}.
    \end{equation*}
\end{lemma}

\begin{proposition}[Lorentz--Strichartz estimates]\label{prel/strichartz}
    Let $\theta_1,\theta_2 \in (0,1)$, $\alpha \in [0,\frac{1}{2}]$, $(q_i,p_i) = (\frac{6}{\theta_i(\alpha + 1)},\frac{2}{1-\theta_i})$ for $i \in \{1,2\}$, and $\frac{1}{p_2} + \frac{1}{p_2'} = \frac{1}{q_2} + \frac{1}{q_2'} = 1$.

    Suppose that $q_1 \leq q_2$ and $q_1 < \infty$. Then for any spacetime slab $J \times \R$ and $t_0 \in J$,
    \begin{equation*}
        \bigg\|\int_{t_0}^t |\partial_x|^\frac{\theta_1\alpha}{2} e^{-(t-s)\partial_x^3} g(s,\cdot) ds\bigg\|_{L_t^{q_1,2}L_x^{p_1}(J)} \lesssim \big\||\partial_x|^{-\frac{\theta_2\alpha}{2}}g\big\|_{L_t^{q_2'}L_x^{p_2'}(J)}.
    \end{equation*}

    Instead, suppose that $q_2 \leq q_1$. Then for any spacetime slab $J \times \R$ and $t_0 \in J$,
    \begin{equation*}
        \bigg\|\int_{t_0}^t |\partial_x|^\frac{\theta_1\alpha}{2} e^{-(t-s)\partial_x^3} g(s,\cdot) ds\bigg\|_{L_t^{q_1,\frac{2q_1}{q_2}}L_x^{p_1}(J)} \lesssim \big\||\partial_x|^{-\frac{\theta_2\alpha}{2}}g\big\|_{L_t^{q_2',2}L_x^{p_2'}(J)}.
    \end{equation*}
\end{proposition}

\begin{proof}
    We mimic the usual proof of non-endpoint Strichartz estimates. Throughout the proof, we estimate over $\R \times \R$; the restriction to the spacetime slab $J \times \R$ can be achieved by replacing $g(s,x)$ by $g(s,x) \mathds{1}_{s \in J}$.

    We first aim to show the $q_2 = \infty, p_2 = 2$ estimate. Arguing by duality, we find
    \begin{align*}
        \bigg\|\int_{t_0}^t |\partial_x|^\frac{\theta_1\alpha}{2} e^{-(t-s)\partial_x^3}&  g(s,\cdot) ds\bigg\|_{L_t^{q_1,2}L_x^{p_1}} \\
        & = \sup\bigg|\bigg\langle\int_{t_0}^t |\partial_x|^\frac{\theta_1\alpha}{2} e^{-(t-s)\partial_x^3} g(s,x) ds,  f(t,x)\bigg\rangle_{t,x}\bigg| \\
        & = \sup\bigg|\int \int_{t_0}^t\Big\langle  e^{s\partial_x^3} g(s,x),|\partial_x|^\frac{\theta_1\alpha}{2}e^{t\partial_x^3}f(t,x)\Big\rangle_{x}dsdt\bigg| \\
        & \leq \sup\big\|e^{s\partial_x^3} g(s,x)\big\|_{L_s^1L_x^2}\bigg\|\int |\partial_x|^\frac{\theta_1\alpha}{2}e^{t\partial_x^3}f(t,x)dt\bigg\|_{L_x^2}.
    \end{align*}
    where the supremum is taken over $\big\{f : \|f\|_{L_t^{q_1',2}L_x^{p_1'}} = 1\big\}$ and $\langle\cdot,\cdot\rangle_*$ is the $L^2$ inner product with subscript indicating the variable(s) over which the inner product is taken.

    By the dual estimate of Lemma \ref{prel/strichartz-basic}, we then find that
    \begin{align*}
        \bigg\|\int_{t_0}^t |\partial_x|^\frac{\theta_1\alpha}{2} e^{-(t-s)\partial_x^3} g(s,\cdot) ds\bigg\|_{L_t^{q_1,2}L_x^{p_1}}
        & \lesssim \sup\|g(s,x)\|_{L_s^1L_x^2}\|f\|_{L_t^{q_1',2}L_x^{p_1'}}\\
        & = \|g(s,x)\|_{L_s^1 L_x^2}.
    \end{align*}
    Interpolating with the estimates from Lemma \ref{prel/strichartz-basic}, this completes the proof of Proposition \ref{prel/strichartz} in the case $q_1 \leq q_2$.

    We now consider $q_1 = \infty$, $p_1 = 2$. Arguing again by duality, we may estimate
    \begin{align*}
        \bigg\|\int_{t_0}^t e^{-(t-s)\partial_x^3} g(s,\cdot) ds\bigg\|_{L_t^\infty L_x^2}
         % & = \sup\bigg|\bigg\langle\int_{t_0}^t e^{-(t-s)\partial_x^3} g(s,x) ds,  f(t,x)\bigg\rangle_{t,x}\bigg| \\
        & = \sup\bigg|\int \int_{t_0}^t\Big\langle  e^{s\partial_x^3} g(s,x),e^{t\partial_x^3}f(t,x)\Big\rangle_{x}dsdt\bigg| \\
        & \leq \sup\big\|e^{t\partial_x^3} f(t,x)\big\|_{L_t^1L_x^2}\bigg\|\int_{t_0}^t e^{s\partial_x^3}g(s,x)ds\bigg\|_{L_x^2} \\
        & \lesssim \Big\||\partial_x|^\frac{-\theta_2\alpha}{2} g(t,x) \Big\|_{L_t^{q_2',2}L_x^{p_2'}}.
    \end{align*}
    where the supremum is taken over $\{f : \|f\|_{L_t^1 L_x^2} = 1\}$. Interpolating with Lemma \ref{prel/strichartz-basic}, this completes the proof of Proposition \ref{prel/strichartz} in the case $q_2 \leq q_1$. Note that the Lorentz exponent $\frac{2q_1}{q_2}$ arises from the observation:
    \begin{equation*}
        L_t^{q_1,\frac{2q_1}{q_2}}\dot{H}_x^{p_1,\frac{\theta_1\alpha}{2}} = \Big( L_t^\infty L_x^2,\;L_t^{q_2,2}\dot{H}_x^{p_2,\frac{\theta_2\alpha}{2}}\Big)_{\left[\frac{q_2}{q_1}\right]},
    \end{equation*}
    in the sense of complex interpolation.
\end{proof}

\end{document}
